All of the numbers in this report come from the two scripts named below.
Run them from their own directory. No argument is needed for the first
one.

\begin{verbatim}
cd cp_matrices/examples_2026/3D_surface

matlab -batch "example_ellipsoid_flipped_cap_branch_consistency_vgmm"

matlab -batch "example_half_twisted_rectangular_tube_rotation_convergence(1./[20 25 30 36 44])"
\end{verbatim}

\subsection{Two changes were needed to produce these runs}

\paragraph{A bug stops the ellipsoid script after its first cut.} In
\texttt{example\_ellipsoid\_flipped\_cap\_branch\_consistency\_vgmm.m},
the local function \texttt{error\_surface\_figure} is declared with
twenty parameters, ending \texttt{geoname, dkappa}. Its body then uses
\texttt{thetaex(1)}:

\begin{verbatim}
                     dkappa, thetaex(1)*180/pi, dx)}, 'FontSize', 9);
\end{verbatim}

\texttt{thetaex} is never passed to it. A local function in a script file
does not share the script's workspace, so this raises
\texttt{Unrecognized function or variable 'thetaex'}. The default
\texttt{z\_3d = 0.390} is not empty, so the figure is always requested,
and the script always aborts after the first cut height. It therefore
never reaches cuts two, three and four, nor the summary tables, nor the
summary figure.

The fix is to pass the variable through. Add \texttt{thetaex} to the end
of the function's parameter list, and to the end of its call. The call
site already has \texttt{thetaex} in scope: it is computed once per cut
height, just above.

\paragraph{The grid spacings were reduced.} The default list is
$h_k = 0.065\times 0.82^{\,k}$ for $k=0,\dots,9$. At four cut heights and
four solves per level, with the band size growing like $h^{-2}$, that is
a run of roughly a day. The runs here use $k=0,\dots,5$, which is the
same list truncated. Nothing else in the script was changed. The
consequences for the fitted rates are discussed in
Section~\ref{sec:scatter}.

\paragraph{Cut one was run separately.} Because the abort above was found
after cut one had completed, the first cut height comes from one process
and the other three from another. The two processes used identical
settings apart from \texttt{s0\_list} and \texttt{z\_3d}. The script's
own end-of-run summary figure therefore shows only three cuts;
Figure~\ref{fig:cutsummary} in this report is drawn from the logged
numbers of both processes and shows all four.

\paragraph{Supporting files.} The two scripts share the rotation
helpers and the geometry routines:

\begin{center}
\begin{tabular}{@{}ll@{}}
  \toprule
  file & what it does \\
  \midrule
  \texttt{rotations/rot2d\_from\_to.m}
    & the plane rotation carrying one direction onto another, by
      Rodrigues' formula \\
  \texttt{rotations/glue\_rot2d.m}
    & the same, with the junction convention: outward onto inward \\
  \texttt{rotations/rotate\_about2d.m}
    & applies a $(\cos,\sin)$ pair about a junction point \\
  \texttt{rotations/rot2d\_err.m}
    & the angle and matrix error of a rotation, through the relative
      rotation \\
  \texttt{3D\_surface/angle3D.m}
    & the same construction in three dimensions \\
  \texttt{surfaces/cpEllipseArc.m}
    & closest point of an arc of an ellipse, with an edge flag \\
  \texttt{surfaces/cpHalfTwistedRectTube.m}
    & closest point of one strip of the tube, with an edge flag \\
  \texttt{surfaces/halfTwistedRectTubeFrame.m}
    & the tube parametrization and its analytic derivatives \\
  \texttt{3D\_surface/halfTwistedRectTubeMMS.m}
    & the manufactured solution of the tube \\
  \texttt{3D\_surface/halfTwistedRectTubeBands.m}
    & the grid, the two bands, and the routing targets, per level \\
  \bottomrule
\end{tabular}
\end{center}

\paragraph{Figures.} Each script writes its own PNG files into
\texttt{examples\_2026/figs/}. The tube script also writes a MAT file
with the full \texttt{results} structure. The two plot-only scripts
\texttt{example\_half\_twisted\_rectangular\_tube\_geometry.m} and
\texttt{example\_half\_twisted\_rectangular\_tube\_narrow\_band.m}
produce Figures~\ref{fig:tubegeom} and~\ref{fig:tubeband}. They solve
nothing.

\paragraph{Theory notes.} \texttt{examples\_2026/angle\_rotation\_glue.md}
carries the derivation. Section~3 covers the glue rotation step by step,
and Section~10 covers the tube.
